\section{Parameters}\label{sec:params}

The parameter sets are produced by \art{crypto:eigen_seta/examples/gen_params.rs},
which writes \art{crypto:eigen_seta/src/params_generated.rs}. They satisfy
four constraints.
\begin{enumerate}
\item[(P1)] $d=3^F\ge 2^{2\lambda}$, with $F=\lceil 2\lambda/\log_2 3\rceil$.
  Message recovery costs about $\sqrt d$ (\S\ref{sec:msg}).
\item[(P2)] $2^E>28\,d^2f_{\max}^2$ with $f_{\max}=2f_{\mathrm{lo}}$. Then
  $m=2^E-7d^2f^2>\tfrac34 2^E$ for every admissible $f$, so the two-squares
  step has a large target. $E$ is the smallest exponent that works.
\item[(P3)] $f_{\mathrm{lo}}=2^\kappa\lfloor p/2^E\rfloor+1$. Then
  $fg\gtrsim 2^\kappa p$, and the norm equation $N(x)+pN(y)=fg$ has about
  $2^\kappa$ candidate values of $y\ne0$. We take $\kappa=12$ for the real sets
  and $\kappa=8$ for the toys.
\item[(P4)] $p=2^{E+1}3^Fc-1$ is prime with $p\equiv3,5,6\pmod7$, and $c$ is the
  smallest cofactor that achieves this.
\end{enumerate}

Since $f\approx 2^{\kappa+1}cd$ and $2^E\approx 28d^2f^2$, we get
$p\approx 2^{E+1}d\approx d^5\cdot 2^{2\kappa+O(1)}c^2$. That is
$\log_2p\approx 10\lambda$. The factor $d^5$ is the price of decrypting with
a single chain, whose torsion order must exceed the degree $7d^2f^2+c^2$ of
$\gamma$. Evaluating $\Phi$ from two half-chains, as in SQIsign~2D-style
implementations~\cite{DMPR23}, would let $2^E$ shrink to about
$\sqrt{\deg\gamma}$. We have not implemented this.

\begin{table}[h]
\centering
\caption[Parameter sets]{Parameter sets. $d=3^F$; ``$f$ bits'' is the bit length of the secret
degree in the recorded keys; public key and ciphertext are four
$\F_{p^2}$ elements each. Source:
\art{crypto:eigen_seta/src/params_generated.rs}; $f$ bits from
\art{crypto:research/eigen_seta_rust_20261011/reports/run-20261011T1813Z-bench/records/bench_output.txt}.}
\label{tab:params}
\begin{tabular}{lrrrrrrrrr}
\toprule
set & $\lambda$ & $\log_2 p$ & $E$ & $F$ & $\log_2 d$ & $c$ & $\kappa$ & $f$ bits & pk = ct (B)\\
\midrule
Toy16 & 16 & 216 & 174 & 21 & 33.3 & 158 & 8 & 51 & 216\\
Toy32 & 32 & 369 & 297 & 41 & 65.0 & 55 & 8 & 80 & 376\\
L1 & 128 & 1341 & 1076 & 162 & 256.8 & 144 & 12 & 278 & 1344\\
L3 & 192 & 1991 & 1594 & 243 & 385.1 & 1059 & 12 & 409 & 1992\\
L5 & 256 & 2630 & 2106 & 324 & 513.5 & 380 & 12 & 536 & 2632\\
\bottomrule
\end{tabular}
\end{table}

The toy sets have no security; they exist so that every code path can be
tested quickly.
